\chapter{Materiais e métodos}
\label{cap:metodos}

Este capítulo descreve como o trabalho foi conduzido: a natureza da pesquisa,
o processo de desenvolvimento, o ambiente e os instrumentos, as métricas
coletadas e suas formas de coleta, os critérios de aceitação e as ameaças à
validade do método. O objetivo é permitir que outro pesquisador reproduza as
medições ou avalie criticamente as decisões tomadas.

\section{Natureza e abordagem da pesquisa}
\label{sec:natureza}

A pesquisa é aplicada e de natureza tecnológica: seu produto é um artefato
--- um jogo sério com IA local --- construído e avaliado tecnicamente. A
abordagem é de estudo de caso único de desenvolvimento
\cite{yin2015casestudy}, com o ciclo projeto--construção--avaliação descrito
na literatura de \emph{design science} \cite{hevner2004design}: o problema e
os requisitos foram fixados primeiro; o artefato foi construído em incrementos
demonstráveis; e cada incremento foi avaliado por evidências definidas \emph{a
priori} (testes automatizados, medições e capturas de tela). As ameaças à
validade e as limitações do desenho estão declaradas na
\autoref{sec:ameacas_metodo}. A formatação deste documento segue as normas
ABNT aplicáveis a trabalhos acadêmicos \cite{nbr14724,nbr10520,nbr6023}; o
modelo institucional definitivo deve ser validado com o autor.

\section{Processo de desenvolvimento}
\label{sec:processo}

O processo adotado é uma adaptação pragmática do desenvolvimento orientado a
especificação, em que a especificação e os critérios de aceitação precedem a
implementação e funcionam como contrato de cada iteração. Três regras
estruturaram o trabalho:

\begin{enumerate}
  \item \textbf{Regra de Ouro:} cada prova de conceito termina em uma versão
  executável e demonstrável, com caminho de teste reproduzível; não se avança
  com o protótipo quebrado.
  \item \textbf{Input semântico:} o \emph{gameplay} consome apenas ações
  nomeadas de um mapa de entrada (por exemplo, mover, interagir, pausar), e
  nunca índices numéricos de botão --- que variam entre adaptadores USB.
  \item \textbf{IA sem controle:} a IA devolve apenas texto e uma intenção de
  um conjunto finito; nenhuma resposta é executada como comando.
\end{enumerate}

A Figura~\ref{fig:processo} resume o ciclo executado. A cada etapa, a
implementação era seguida de testes automatizados; somente com a suíte verde o
projeto avançava. Evidências visuais (capturas) e medições eram registradas
junto com a atualização dos documentos de estado do projeto
\cite{lab404repo}.

\begin{figure}[htbp]
\centering
\includegraphics[width=0.85\textwidth]{05_processo.png}
\caption{Ciclo de desenvolvimento incremental adotado: especificação, implementação, verificação automatizada e registro de evidências.}
\label{fig:processo}
\end{figure}

O desenvolvimento foi assistido por agentes de inteligência artificial no
editor de código. Para que essa assistência permanecesse sob controle, os
documentos de especificação, regras e estado foram mantidos versionados e
foram tratados como a fonte de intenção: cada tarefa delegada a um agente
referenciava requisitos e critérios de aceitação explícitos. Um servidor MCP
\cite{mcp2026} expôs a cena do Blender à manipulação programática, o que
permitiu conduzir a modelagem 3D e a exportação também dentro do fluxo
assistido. A Figura~\ref{fig:arquitetura} (no
\autoref{cap:desenvolvimento}) mostra o resultado arquitetural desse processo.

\section{Ambiente de referência}
\label{sec:ambiente}

Todas as medições deste trabalho foram realizadas em uma única máquina de
referência, descrita na Tabela~\ref{tab:ambiente}. A escolha é deliberada:
o critério de projeto é a execução em computador pessoal sem GPU dedicada, e
por isso nenhuma medição utiliza hardware especializado.

% Tabela de ambiente de referência (usada no capítulo 4).
\begin{table}[htbp]
\centering
\caption{Ambiente de referência utilizado em todas as medições (verificado em 2026-09-10).}
\label{tab:ambiente}
\footnotesize
\begin{tabular}{@{}>{\raggedright\arraybackslash}p{4.4cm} >{\raggedright\arraybackslash}p{9.4cm}@{}}
\toprule
Item & Versão ou configuração \\
\midrule
Sistema operacional & Ubuntu 24.04.4 LTS (64 bits) \\
Processador & AMD Ryzen 7 5700U com Radeon Graphics \\
Memória principal & 30 GiB \\
GPU & GPU integrada AMD Radeon Graphics (sem GPU dedicada) \\
Motor de jogo & Godot 4.5 (\texttt{4.5.stable.official.876b29033}), renderizador \texttt{gl\_compatibility} \\
Modelagem 3D & Blender 5.2.1 LTS \\
Linguagem do adaptador & Python 3.12.3 \\
Bibliotecas do adaptador & FastAPI 0.141.1; Uvicorn 0.52.4; Pydantic 2.13.5; requests 2.34.2 \\
Inferência local & Ollama; modelo-alvo \texttt{llama3.2:3b} (não instalado no ambiente); validação realizada com \texttt{llama3.1:8b} \\
Formato de assets & glTF 2.0 binário (GLB) \\
Resolução de execução & 1280$\times$720 com sincronismo vertical (vsync) \\
\bottomrule
\end{tabular}
\par\vspace{2pt}{\footnotesize Fonte: elaborado pelo autor (2026); valores obtidos por comando na máquina de referência.}
\end{table}


O projeto mantém um passo de importação de recursos (modelos GLB e áudio)
executado por linha de comando, necessário na primeira execução de cada
cópia do repositório; a suíte de testes e as capturas também são executadas
por linha de comando, sem interface gráfica (quando aplicável) nem
intervenção manual. A descrição do repositório, dos comandos e da organização
dos diretórios está no \autoref{ap:repositorio}.

\section{Instrumentos e métricas}
\label{sec:instrumentos}

A Tabela de instrumentos resume as métricas coletadas, o instrumento
correspondente e a forma de coleta. Todas as métricas são produzidas por
execução real do sistema na máquina de referência; nenhuma é estimada.

\begin{quadro}[htbp]
\centering
\caption{Instrumentos de avaliação e métricas coletadas.}
\label{quad:instrumentos}
\footnotesize
\begin{tabular}{@{}>{\raggedright\arraybackslash}p{3.6cm} >{\raggedright\arraybackslash}p{4.2cm} >{\raggedright\arraybackslash}p{5.4cm}@{}}
\toprule
Métrica & Instrumento & Forma de coleta \\
\midrule
Verificações funcionais por etapa & Cenas de teste \emph{headless} do motor (uma por POC) & Execução de \texttt{tests/run\_all.sh}; contagem de verificações aprovadas impressa por cena \\
Testes de contrato da IA & Suíte do adaptador em Python & Execução por linha de comando contra o serviço local \\
Desempenho gráfico & Motor em execução com janela & Leitura de quadros por segundo do próprio motor, com sincronismo vertical, em 1280$\times$720 \\
Latência da conversa & Cronometragem de três execuções independentes & Chamada direta ao adaptador (partida fria), terminal no jogo e cliente em teste ao vivo; valores obtidos do relatório de execução e da captura de evidência \\
Complexidade do asset 3D & Inspeção programática do arquivo GLB & Contagem de nós, materiais e malhas de colisão por script \\
Evidências visuais & Arnês de captura do projeto & Execução com variáveis de ambiente que posicionam o jogador, opcionalmente concluem a missão e salvam imagens em 1280$\times$720 \\
\bottomrule
\end{tabular}
\par\vspace{2pt}{\footnotesize Fonte: elaborado pelo autor (2026).}
\end{quadro}

\section{Critérios de aceitação e rastreabilidade}
\label{sec:criterios}

Cada requisito funcional possui um critério de aceitação no formato
DADO\slash QUANDO\slash ENTÃO, escrito antes da implementação. A
Tabela~\ref{tab:aceitacao_resumo} (no \autoref{cap:resultados}) resume o
status de cada critério; o Quadro~\ref{quad:ca_exemplos} apresenta dois
exemplos representativos, um de cada subsistema crítico. A lista completa dos
requisitos e critérios está no \autoref{ap:requisitos}.

\begin{quadro}[htbp]
\centering
\caption{Exemplos de critérios de aceitação usados como contrato de implementação.}
\label{quad:ca_exemplos}
\footnotesize
\begin{tabular}{@{}>{\raggedright\arraybackslash}p{1.4cm} >{\raggedright\arraybackslash}p{11.6cm}@{}}
\toprule
ID & Critério (resumo) \\
\midrule
CA-009 & DADO que o serviço de IA está indisponível ou lento, QUANDO o jogador envia uma mensagem, ENTÃO o jogo exibe uma resposta de \emph{fallback} dentro do limite de tempo, sem travar. \\
CA-010 & DADO uma resposta do modelo, QUANDO a intenção não pertence ao conjunto permitido, ENTÃO ela é convertida para \texttt{NONE} e a resposta textual permanece válida. \\
\bottomrule
\end{tabular}
\par\vspace{2pt}{\footnotesize Fonte: elaborado pelo autor (2026), a partir da especificação do projeto.}
\end{quadro}

A rastreabilidade é mantida em duas matrizes. A primeira associa cada
objetivo específico desta monografia aos métodos e evidências correspondentes;
a segunda associa requisito, critério de aceitação, teste e evidência. Ambas
estão consolidadas no arquivo de rastreabilidade do projeto
(MONOGRAFIA\_RASTREABILIDADE.md) e resumidas no
\autoref{cap:resultados}.

\section{Procedimentos experimentais}
\label{sec:procedimentos}

Os procedimentos seguiram uma ordem fixa, executada ao final do
desenvolvimento e repetida a cada alteração relevante:

\begin{enumerate}
  \item \textbf{Verificação funcional:} execução completa da suíte
  \emph{headless} (etapas \poc{1} a \poc{6}) e da suíte do adaptador de IA,
  registrando o número de verificações e falhas de cada cena;
  \item \textbf{Verificação de integração ao vivo:} execução do jogo com o
  adaptador e o modelo local ativos, enviando uma pergunta real pelo terminal
  da IA e registrando a resposta e a latência;
  \item \textbf{Medição de latência:} três contextos independentes --- chamada
  direta ao adaptador em partida fria, mensagem pelo terminal com o jogo em
  execução e teste ao vivo do cliente ---, todos com o mesmo modelo local;
  \item \textbf{Medição de desempenho:} execução da cena completa em
  1280$\times$720 com sincronismo vertical, lendo os quadros por segundo
  reportados pelo motor;
  \item \textbf{Evidências visuais:} geração de capturas em poses
  pré-determinadas, incluindo o estado inicial, o estado após a missão
  concluída e o terminal da IA com a resposta do modelo;
  \item \textbf{Inspeção do artefato 3D:} medição programática do GLB
  exportado (nós, materiais, colisões e tamanho em disco).
\end{enumerate}

Os itens que exigem presença física ou julgamento humano --- operação de
mouse com captura de cursor, calibração do adaptador de joystick em hardware
real e o percurso completo do início ao fim por uma pessoa --- estão
declarados como pendentes de validação manual e não foram contabilizados como
concluídos.

\section{Ameaças à validade do método}
\label{sec:ameacas_metodo}

Quatro ameaças devem ser consideradas na leitura dos resultados. A primeira é
de \textbf{generalização}: trata-se de um estudo de caso único, com um
artefato construído por um único autor; os resultados não se generalizam
automaticamente para outros jogos, outros modelos de linguagem ou outras
máquinas. A segunda é de \textbf{medição}: as latências foram coletadas em
três execuções por contexto, sem repetição estatística; são indicativas, não
experimentais. A terceira é de \textbf{autoria}: o mesmo autor especificou,
implementou e avaliou o artefato, o que aumenta o risco de viés de confirmação;
a mitigação adotada é externa ao autor --- critérios escritos antes da
implementação, testes automatizados executáveis por terceiros e artefatos
públicos. A quarta é de \textbf{ambiente}: os resultados de desempenho valem
para a máquina de referência; o capítulo de discussão registra duas
particularidades de ambiente que afetaram a experiência de uso e que outros
projetos podem encontrar (confinamento de pacote do motor e estado de áudio
por aplicativo).

\section{Planejado versus executado}
\label{sec:planejado_executado}

O trabalho distingue explicitamente o que foi planejado do que foi executado.
O Quadro~\ref{quad:planejado} apresenta a distinção por componente; a coluna
de status registra o que permanece aberto, sem conversão de intenção em
resultado.

\begin{quadro}[htbp]
\centering
\caption{Planejado versus executado, por componente do projeto.}
\label{quad:planejado}
\footnotesize
\begin{tabular}{@{}>{\raggedright\arraybackslash}p{4.0cm} >{\raggedright\arraybackslash}p{4.6cm} >{\raggedright\arraybackslash}p{4.2cm}@{}}
\toprule
Componente & Planejado & Executado e status \\
\midrule
Jogo (etapas \poc{1}--\poc{6}) & Seis provas de conceito demonstráveis & Executado; suíte automatizada verde \\
IA local & Conversa com modelo local e \emph{fallback} & Executado com \texttt{llama3.1:8b}; modelo-alvo \texttt{llama3.2:3b} não instalado (aberto) \\
Desempenho & Executar sem GPU dedicada & Executado a 60 FPS (vsync) na máquina de referência \\
Medições de latência & Três contextos independentes & Executado (uma execução por contexto) \\
Evidências visuais & Capturas em poses controladas e vídeo & Executado (capturas e vídeo de 68 s) \\
Validação de hardware & Operar com joystick físico e mouse & Pendente (validação manual) \\
Percurso humano (CA-013) & Jogo completo por uma pessoa & Pendente (validação manual) \\
Empacotamento (FR-033) & Executável independente & Bloqueado por ausência de modelos de exportação no ambiente \\
Estudo com usuários & Fora do escopo declarado & Não realizado (trabalho futuro) \\
\bottomrule
\end{tabular}
\par\vspace{2pt}{\footnotesize Fonte: elaborado pelo autor (2026).}
\end{quadro}

\section{Considerações do capítulo}
\label{sec:consideracoes_metodos}

O método combina elementos de engenharia de software --- especificação,
rastreabilidade, testes automatizados --- com os procedimentos de medição
descritos nesta seção. É esse conjunto que sustenta a afirmação, repetida ao
longo do trabalho, de que cada número apresentado possui origem verificável:
ele foi produzido por um comando, por uma execução ou por uma medição, e não
por estimativa.
